\subsection{ADO 定理}

命题 1. 设 \(\mathfrak{g}\) 是李代数，\(\mathfrak{n}\) 是其最大幂零理想，\(\mathfrak{a}\) 是 \(\mathfrak{g}\) 的幂零理想，\(\rho\) 是 \(\mathfrak{a}\) 的有限维表示，且 \(\rho(a)\) 中每个元素都幂零。则 \(\rho\) 有一个有限维扩张 \(\sigma\) 到 \(\mathfrak{g}\)，使 \(\sigma(n)\) 中每个元素都幂零。

令 \(\mathfrak{a} = \mathfrak{n}_0 \subset \mathfrak{n}_1 \subset \cdots \subset \mathfrak{n}_p = \mathfrak{n}\) 是 \(\mathfrak{n}\) 的子代数序列，其中 \(\mathfrak{n}_{i-1}\) 是 \(\mathfrak{n}_i\) 的理想，且当 \(1 \leqslant i \leqslant p\) 时余维为 1。因此 \(\mathfrak{n}_i\) 是 \(\mathfrak{n}_{i-1}\) 与一个一维子代数的直和。由于 \(\operatorname{ad}_x\) 对所有 \(x \in \mathfrak{n}\) 都幂零，可以（定理 1）逐个找到 \(\rho_1, \rho_2, \ldots, \rho_p = \rho'\)，它们是 \(\rho\) 到 \(\mathfrak{n}_1, \mathfrak{n}_2, \ldots, \mathfrak{n}_p = \mathfrak{n}\) 的有限维扩张，使 \(\rho'(\mathfrak{n})\) 中每个元素幂零。

令 \(\mathfrak{r}\) 为 \(\mathfrak{g}\) 的根基，并令 \(\mathfrak{n} = \mathfrak{r}_0 \subset \mathfrak{r}_1 \subset \dots \subset \mathfrak{r}_q = \mathfrak{r}\) 为 \(\mathfrak{r}\) 的子代数序列，其中 \(\mathfrak{r}_{i-1}\) 是 \(\mathfrak{r}_i\) 的理想，且当 \(1 \leqslant i \leqslant q\) 时余维为 1。因此 \(\mathfrak{r}_i\) 是 \(\mathfrak{r}_{i-1}\) 与一个一维子代数的直和。由于 \([\mathfrak{r}, \mathfrak{r}] \subset \mathfrak{n}\)，可以（定理 1）逐个找到 \(\rho_1', \rho_2', \ldots, \rho_q' = \rho''\)，它们是 \(\rho'\) 到 \(\mathfrak{r}_1, \mathfrak{r}_2, \ldots, \mathfrak{r}_q = \mathfrak{r}\) 的有限维扩张，使 \(\rho''(\mathfrak{n})\) 中每个元素幂零。

最后 g 是 r 与一个子代数 s 的直和（§ 6，第 8 号，定理 5）。由于 \([s, r] \subset n\)，可以（定理 1）找到一个从 ρ'' 到 g 的有限维扩张 σ，使 σ(n) 中每个元素幂零。

定理 2. 每个李代数都有有限维忠实线性表示。

更精确地说：

定理 3（Ado）。设 \(\mathfrak{g}\) 是李代数，\(\mathfrak{n}\) 是其最大幂零理想。存在有限维忠实表示 \(\rho\)（它是 \(\mathfrak{g}\) 的表示），使 \(\rho(\mathfrak{n})\) 中每个元素幂零。

一维李代数 K 有限维忠实表示 \(\tau\)，使 \(\tau(\mathbf{K})\) 中每个元素幂零，例如表示

\[
\lambda \mapsto \left( \begin{array}{c c} 0 & 0 \\ \lambda & 0 \end{array} \right).
\]

容易推出，g 的中心 \(c\)（即 \(g\) 的中心）是若干一维代数的乘积，因此存在其有限维忠实表示 \(\sigma\)，使 \(\sigma(c)\) 中每个元素幂零。令 \(\sigma_1\) 为 \(\sigma\) 向 \(g\) 的有限维扩张，使 \(\sigma_1(n)\) 中每个元素幂零（命题 1）；若 \(t\) 表示 \(\sigma_1\) 的核，则 \(t \cap c = \{0\}\)。另一方面，令 \(\sigma_2\) 为 \(g\) 的伴随表示，其核是 \(c\)；\(\sigma_2(n)\) 中每个元素幂零。令 \(\rho\) 为 \(\sigma_1\) 与 \(\sigma_2\) 的直和，它是有限维的；\(\rho(n)\) 中每个元素幂零；且 \(\rho\) 的核包含于 \(t\) 和 \(c\) 中，因而为零，所以 \(\rho\) 是忠实的。

